\chapter{Introduction}

Molecular property prediction refers to the computational estimation of chemical, physical, or biological characteristics of molecules based on their structure. These properties—ranging, e.g., from solubility and toxicity to binding affinity and biological activity—are essential for understanding molecular behavior in a wide array of contexts. Accurate predictions promise to significantly accelerate early-stage drug discovery, reduce the need for costly and time-consuming laboratory experiments, and support advances in molecular medicine, materials science and drug discovery.
The reliable prediction of molecular properties is of great interest for the understanding of molecular mechanisms in biology and crucial for molecular medicine as well as drug discovery. Computational  methods for this task rely on processable representations of the molecular structure.
To enable such predictions, molecules must be translated into formats that computational models can interpret and process. A variety of molecular representations has been developed to capture the relevant structural and chemical information. Among the most widely used are circular molecular fingerprints, such as the extended-connectivity fingerprints (ECFPs), which encode local atomic neighborhoods into fixed-length vectors. These fingerprints have demonstrated strong performance when paired with traditional machine learning models, particularly in tasks like quantitative structure–activity relationship (QSAR) modeling.
A variety of molecular representations have been developed to capture the relevant structural and chemical information. Among the most widely used are circular molecular fingerprints, such as the extended-connectivity fingerprints (ECFPs), which encode local atomic neighborhoods into fixed-length vectors. 

In recent years, the advent of Graph Neural Networks (GNNs), which treat molecules as graphs, where atoms are represented by nodes and bonds by edges, has introduced a new paradigm for molecular representation learning.  By operating directly on molecular graphs, GNNs aim to learn task-specific embeddings through message-passing mechanisms that are conceptually similar to the iterative procedures of fingerprint generation. 
This flexible data-driven approach, as well as the promising performance of neural networks in many other tasks, has led to the expectation that GNNs may outperform handcrafted fingerprints in predictive performance. Their graph-based architecture naturally aligns with the inherent topology of molecules, and their message-passing mechanisms are conceptually similar to the iterative processes used in fingerprint generation, but allow for learnable embeddings, instead of the static, fixed-size nature of the fingerprints.
However, recent empirical studies challenge this assumption, showing that GNNs often fail to consistently surpass classical approaches, particularly in quantitative structure–activity relationship (QSAR) modeling and related tasks (examples for reviews are given in the next chapters).

Despite these theoretical advantages, recent studies have shown that GNNs do not consistently outperform traditional fingerprint-based methods. In some cases, particularly for specific prediction tasks or under data constraints, they even lag behind simpler models in both accuracy and generalizability. This counterintuitive observation has sparked renewed interest in understanding the underlying reasons for this performance gap and in exploring the relationship between learned and handcrafted molecular representations.
Despite the theoretical advantages of GNNs, recent studies have revealed that they do not always outperform traditional fingerprint-based methods. In some cases, particularly for specific tasks such as classification of drug-target interactions or regression of solubility, GNNs have been shown to underperform compared to simpler models using molecular fingerprints. This unexpected result suggests that the performance of different models can vary depending on the nature of the task, the available data, and the design of the model itself.

This discrepancy raises fundamental questions about the nature and limits of learned molecular representations. 
It prompts a deeper investigation into the relative strengths of GNNs versus traditional fingerprints, their respective inductive biases, and the influence of factors such as dataset characteristics, model architecture, and training data availability.

This thesis aims to investigate and compare the performance of Graph Neural Networks and classical molecular fingerprints in a variety of molecular property prediction tasks, including classification as well as regression. By systematically comparing these methods across different datasets and tasks, it seeks to clarify the reasons behind the performance differences and explore how factors such as model architecture, dataset size, and task complexity influence the predictive performance of these approaches.

Circular chemical fingerprints, like the extended connectivity fingerprint (ECFP) based on the Morgan algorithm which iteratively evaluates the neighborhood of each atom to generate a fixed-size vector representation, are commonly used for this purpose\cite{Rogers.2010}. The introduction of graph neural networks 
(GNN) in cheminformatics and molecular biology aimed at generalizing these chemical 
fingerprints\cite{Duvenaud.30.09.2015,Kearnes.2016}. GNNs seem to be particularly suited for this task since they can directly work on the molecular graph representation and graph convolutional layers perform analogous operations to the iterative descriptions of ECFPs. 
However, despite the higher flexibility of neural network representations due to learnable weights, several recent studies have shown that GNN representations do not generally outperform circular fingerprints coupled to a multi-layer perceptron in quantitative structure–activity relationship (QSAR) modelling \cite{Deng.2023,Menke.2021,Sabando.2021}. Depending on the specific task, e.g. the prediction of drug targets \cite{Stepisnik.2021}, solubility \cite{Mayr.2018} or the classification of activity cliffs \cite{Jiang.2021}, i.e. pairs of small molecules with high structural similarity, but also unexpectedly high difference in binding affinity regarding a specific target, they even perform significantly worse on some benchmarks. Similarly, molecular presentations of state-of-the-art GNNs do not outperform prior approaches in protein-ligand-interaction prediction \cite{Volkov.2022}. This behaviour calls for an investigation of possible reasons, e.g. overfitting due to the higher expressivity and biased datasets, and of its dependence on specific parameters, e.g. the amount of training data. 

In particular, it is envisaged to elucidate the capricious performance of graph neural 
networks in molecular property prediction tasks and explore their connection to classical chemical fingerprints as well as the reasons for the difference in performance through a detailed analysis of both methods. 
A central part will consist in the reimplementation of a fingerprint as graph neural network. In the context of an ablation study, neural layers will be stepwise added to this basic representation and the effect of the involved operations, e.g. convolutions, will be systematically explored. The impact of different layers and techniques, e.g. dropout mechanisms, can be studied as well as the effect of hyperparameters. Finally, these computer experiments should contribute to an explanation of the performance and limits of graph neural networks in molecular property prediction and similar tasks. 

